\section{The missing conversion from curve degree to matrix degree}
Give $R=\C[x,y,z,w]/(xw-yz-1)$ its minimum representative total-degree
filtration. For a finite right subgroup $\Gamma\subset SL_2$ give
$A=R^\Gamma$ the induced filtration.
\begin{lemma}
The associated graded rings of $R$ and $A$ are domains. Consequently, for
nonzero $f,g\in A$, $\deg(fg)=\deg f+\deg g$. If $h$ is nonconstant and
$p\in\C[T]$ has degree $d>0$, then $\deg p(h)=d\deg h$.
\end{lemma}
\begin{proof}
The highest part of $(xw-yz-1)b$ is $(xw-yz)b_{\rm top}$. Thus
$\operatorname{gr}R=\C[x,y,z,w]/(xw-yz)$, a domain. The group action preserves
the filtration. Averaging over the finite group makes invariants exact
on each filtration quotient, so
$\operatorname{gr}(R^\Gamma)=(\operatorname{gr}R)^\Gamma$, also a domain.
The highest product symbol is nonzero. In $p(h)$ its top power has
strictly greater degree than every other power and cannot cancel.
\end{proof}

The existing normalization theorem in
\texttt{COMPLEX\_HEIGHT\_IMAGES\_AND\_CAUCHY\_RIGIDITY.md} says that a
nonconstant phase on either full exceptional quotient with one-dimensional
compact image is $P=p(h)$ for a regular real-valued seed $h$. Its finite
curve-marker theorem says that a nonzero pure such phase has an infinite
nonzero marked tail for one fixed $h^j$, $1\le j<d=\deg p$.
The existing Cauchy-transform exclusion says its image $p(h(M))$ must
separate the complex plane.

\begin{proposition}
Let $P\ne0$ be pure-null on $SL_2/2O$ or $SL_2/2I$, and suppose its compact
image has real dimension at most one. Put $D=\deg P$ in matrix entries.
Then $D\ge24$ for $2O$, and $D\ge36$ for $2I$. A fixed bad marker $Q$ exists
with $\deg Q\le D-8$, respectively $\deg Q\le D-12$.
At the first possible degree the normalized polynomial $p$ is cubic and
$\deg h$ is exactly $8$, respectively $12$; a bad marker is $h$ or $h^2$.
\end{proposition}
\begin{proof}
Write $P=p(h)$ and $H=\deg h$. The existing exact character/Molien input
GEXCEPTIONALDEGREES20261115 gives $H\ge8$, respectively $H\ge12$.
A polynomial of degree at most two maps a real interval to a segment or
an injective arc: for a quadratic,
$p(t)-p(s)=(t-s)(a(t+s)+b)$. If $a,b$ are real-linearly independent it is
injective; otherwise its image lies in an affine real line. Such an image
does not separate the plane. The retained Cauchy exclusion therefore
forces $d\ge3$. The lemma gives $D=dH\ge3H$. For the existing fixed bad
marker $h^j$, $\deg h^j=jH\le(d-1)H=D-H$. Equality of $D$ with $24$ or $36$
forces $d=3$ and the minimum $H$.
\end{proof}
These bounds concern thin images of the PHASE, not of its auxiliary Hopf
height, and do not exclude full-image counterexamples in smaller degrees.
The thresholds are necessary, not existence statements.

